% =====================================================================
% Mục 3. Tiền xử lý số liệu (bản rút gọn, dùng thay cho 03_preprocessing.tex)
%
% Gói cần có: booktabs, tabularx, array, listings, xcolor, tikz (positioning, arrows.meta),
%   caption, float. Trước khi \input file này phải nạp
%   % Generated by scripts/reports/report_preprocessing.py; do not edit by hand.
\newcommand{\PreRowsRaw}{3406}
\newcommand{\PreColsRaw}{34}
\newcommand{\PreColsSelected}{18}
\newcommand{\PreCellsMarkedMissing}{6474}
\newcommand{\PreReleaseYearMin}{2003}
\newcommand{\PreReleaseYearMax}{2017}
\newcommand{\PreCoreSpeedDash}{907}
\newcommand{\PreCoreSpeedBlank}{29}
\newcommand{\PreUnknownReleaseDate}{30}
\newcommand{\PreBandwidthInMb}{4}
\newcommand{\PreRopsMultiplier}{305}
\newcommand{\PreTextureUntestable}{31}
\newcommand{\PreBandwidthMismatchNotRepaired}{7}
\newcommand{\PreBandwidthMismatchShifted}{3}
\newcommand{\PreTextureMismatchBefore}{274}
\newcommand{\PreTextureMismatchAfter}{0}
\newcommand{\PreBandwidthMismatchBefore}{312}
\newcommand{\PreBandwidthMismatchAfter}{4}
\newcommand{\PreEncodingFixed}{1}
\newcommand{\PreRepairedYears}{2016--2017}
\newcommand{\PreRepairedInFinal}{100}
\newcommand{\PreNotebookYes}{709}
\newcommand{\PreNotebookNo}{2697}
\newcommand{\PreManufacturerLevels}{4}
\newcommand{\PreManufacturerNames}{Nvidia, AMD, Intel, ATI}
\newcommand{\PreMemoryTypeLevels}{12}
\newcommand{\PreMemoryTypeMissing}{56}
\newcommand{\PreMemoryTypeLevelsFinal}{10}
\newcommand{\PreMemoryTypeGroups}{6}
\newcommand{\PreMemoryTypeOtherNames}{DDR, GDDR2, GDDR4, HBM-1, HBM-2}
\newcommand{\PreIntegrated}{374}
\newcommand{\PreUnknownDedicated}{14}
\newcommand{\PreRowsDedicated}{3018}
\newcommand{\PreRowsSingleGpu}{2775}
\newcommand{\PreRowsReleased}{2762}
\newcommand{\PreRowsScoped}{2762}
\newcommand{\PreRowsFinal}{2057}
\newcommand{\PreIncompleteRemoved}{705}
\newcommand{\PreIntelFinal}{0}
\newcommand{\PreMultiGpuRemoved}{243}
\newcommand{\PreUnreleasedRemoved}{13}
\newcommand{\PreDuplicatesRemoved}{0}
\newcommand{\PreLaptopAuxPower}{37}
\newcommand{\PreColsFinal}{16}
\newcommand{\PreDesktop}{1786}
\newcommand{\PreLaptop}{271}
\newcommand{\PreInvalidPowerRows}{1}
\newcommand{\PreInvalidPowerSiblings}{2}
\newcommand{\PreInvalidPowerSiblingWatts}{127}
\newcommand{\PrePowerMissing}{482}
\newcommand{\PrePowerMissingPct}{17,5}
\newcommand{\PreMissingMeanYear}{2009,0}
\newcommand{\PreObservedMeanYear}{2012,7}
\newcommand{\PreMissingYearZ}{-19{,}20}
\newcommand{\PreMissingYearP}{$<$ 0,001}
\newcommand{\PreMissingAmdPct}{53,2}
\newcommand{\PreObservedAmdPct}{38,9}
\newcommand{\PreMissingMakerZ}{5{,}71}
\newcommand{\PreMissingMakerP}{$<$ 0,001}
\newcommand{\PreAtiCards}{83}
\newcommand{\PreAtiMissing}{0}
\newcommand{\PreIntelCards}{7}
\newcommand{\PreLostPowerMean}{63,8}
\newcommand{\PreLostPowerZ}{-14{,}31}
\newcommand{\PreLostPowerP}{$<$ 0,001}
\newcommand{\PreLostMeanYear}{2008,0}
\newcommand{\PreKeptMeanYear}{2013,1}
\newcommand{\PreFinalPowerMean}{124,9}
\newcommand{\PreFinalPowerMedian}{120}
\newcommand{\PreAllPowerMean}{119,0}
\newcommand{\PreAllPowerMedian}{110}
\newcommand{\PreFinalDesktopAuxPct}{71,1}
\newcommand{\PreAllDesktopAuxPct}{68,5}
\newcommand{\PreUnknownYearRows}{30}
\newcommand{\PreUnknownYearPowerMissingPct}{6,7}
\newcommand{\PreUnknownYearDroppedPct}{100,0}
\newcommand{\PreCoreExplainedPct}{60}
\newcommand{\PreCoreMissingInFinal}{220}
\newcommand{\PreCoreMissingInFinalPct}{10,7}
\newcommand{\PreCoreKnownInFinal}{1837}
\newcommand{\PreRowsIfCoreRequired}{1837}
\newcommand{\PreCoreProcessCorr}{-0{,}74}
\newcommand{\PreCoreYearCorr}{0{,}75}
\newcommand{\PreCoreWithinChipSpreadPct}{6}
\newcommand{\PreCoreBetweenChipSpreadPct}{26}
\newcommand{\PreCoreWithinChipElasticity}{0,5}
\newcommand{\PreCoreWithinChipCards}{1627}
\newcommand{\PreCoreWithinChipChips}{132}
\newcommand{\PreCoreMissingPowerMean}{99,7}
\newcommand{\PreCoreKnownPowerMean}{127,9}
\newcommand{\PreNvidiaFinalMean}{126,7}
\newcommand{\PreNvidiaAllMean}{118,7}
\newcommand{\PreNvidiaMinusAmdFinal}{4{,}1}
\newcommand{\PreNvidiaMinusAmdAll}{-0{,}4}
\newcommand{\PreMemoryTypeOtherSize}{18}
\newcommand{\PrePowerFarFromChipMedian}{11}
\newcommand{\PreLowPowerLaptopCards}{11}
\newcommand{\PreRowsBeforeTwoThousandSix}{7}
\newcommand{\PreRowsWithPower}{2280}
\newcommand{\PreRowsIncomplete}{705}
\newcommand{\PreRowsIncompletePct}{25,5}
\newcommand{\PrePowerRowsLostByDropAll}{223}
\newcommand{\PrePowerMeanDropped}{119,0}
\newcommand{\PrePowerMeanImputed}{117,4}
\newcommand{\PrePowerSdDropped}{77,5}
\newcommand{\PrePowerSdImputed}{70,5}
\newcommand{\PrePowerCiDropped}{3,2}
\newcommand{\PrePowerCiImputed}{2,6}
\newcommand{\PrePowerAtMedianDropped}{1,0}
\newcommand{\PrePowerAtMedianImputed}{18,3}
\newcommand{\PreMissingPctGpuId}{0,0}
\newcommand{\PreMissingPctName}{0,0}
\newcommand{\PreMissingPctManufacturer}{0,0}
\newcommand{\PreMissingPctSegment}{0,0}
\newcommand{\PreMissingPctMemoryType}{0,7}
\newcommand{\PreMissingPctMemoryTypeGroup}{0,7}
\newcommand{\PreMissingPctArchitecture}{1,5}
\newcommand{\PreMissingPctReleaseYear}{1,1}
\newcommand{\PreMissingPctMaxPowerW}{17,5}
\newcommand{\PreMissingPctNeedsAuxPower}{17,5}
\newcommand{\PreMissingPctTextureRatePerWatt}{24,2}
\newcommand{\PreMissingPctTmus}{13,0}
\newcommand{\PreMissingPctMemoryBandwidthGbS}{1,3}
\newcommand{\PreMissingPctMemoryMb}{0,9}
\newcommand{\PreMissingPctProcessNm}{10,5}
\newcommand{\PreMissingPctCoreSpeedMhz}{22,0}

% (sinh bằng: python scripts/reports/report_preprocessing.py; script này cũng sinh
% generated/preprocessing_missing_by_period.tex và
% generated/preprocessing_power_by_group.tex). Mọi con số tính từ dữ liệu là macro
% \Pre..., nên chạy lại code là báo cáo tự cập nhật. Giữ nguyên các nhãn của bản đầy
% đủ (sec:preprocessing, sec:missing, tab:rows, tab:columns, ...).
% =====================================================================

\section{Tiền xử lý số liệu}
\label{sec:preprocessing}

Phần này chuyển tệp \texttt{All\_GPUs.csv} (\PreRowsRaw{} dòng, \PreColsRaw{} biến) thành
bộ dữ liệu cho các phân tích ở mục sau. Quy trình viết bằng Python (\texttt{pandas}),
chạy lại từ đầu cho cùng một kết quả và gồm sáu bước (Hình~\ref{fig:preprocess}). Nhóm
theo ba nguyên tắc: không sửa tay tệp gốc; hàm đọc mỗi biến số chỉ nhận đúng định dạng đã
biết và dừng chương trình nếu gặp giá trị lạ; không điền giá trị khuyết. Ngoại lai và
biến đổi logarit được xử lý trong từng phân tích ở các mục sau.

\begin{figure}[!htb]
  \centering
  \begin{tikzpicture}[
      node distance=2.5mm,
      box/.style={rectangle, rounded corners, draw, fill=blue!6, text width=12cm,
        align=left, minimum height=6mm, font=\footnotesize},
      io/.style={rectangle, draw, fill=orange!12, text width=12cm, align=center,
        minimum height=6mm, font=\footnotesize},
      arrow/.style={-Stealth, thick}]
    \node[io] (raw) {Tệp gốc: \PreRowsRaw{} dòng $\times$ \PreColsRaw{} biến};
    \node[box, below=of raw] (read) {1. Đọc mọi cột dưới dạng chuỗi
      (\texttt{read\_raw})};
    \node[box, below=of read] (select) {2. Chọn \PreColsSelected{} biến
      (\texttt{select\_columns})};
    \node[box, below=of select] (clean) {3. Xóa khoảng trắng, đánh dấu ô khuyết, sửa
      lỗi mã hóa (\texttt{clean\_text})};
    \node[box, below=of clean] (parse) {4. Tách đơn vị, đổi kiểu, sửa lỗi ``(x$N$)'',
      bỏ giá trị sai (\texttt{parse\_columns}, \texttt{check\_consistency})};
    \node[box, below=of parse] (filter) {5. Lọc: GPU rời, một GPU, đã phát hành, không
      trùng, đủ biến bắt buộc (\texttt{filter\_rows})};
    \node[box, below=of filter] (derive) {6. Tạo biến mới, giữ các cột cần dùng
      (\texttt{derive\_variables})};
    \node[io, below=of derive] (final) {Dữ liệu cuối: \PreRowsFinal{} dòng $\times$
      \PreColsFinal{} cột};
    \foreach \a/\b in {raw/read, read/select, select/clean, clean/parse,
        parse/filter, filter/derive, derive/final}
      \draw[arrow] (\a) -- (\b);
  \end{tikzpicture}
  \caption{Quy trình tiền xử lý (tên hàm Python trong ngoặc)}
  \label{fig:preprocess}
\end{figure}

% ---------------------------------------------------------------------
\subsection{Chọn biến và làm sạch chuỗi (bước 1--3)}

Mọi cột được đọc dưới dạng chuỗi để từng giá trị đều đi qua bước kiểm tra định dạng; mỗi
dòng được gán mã \texttt{gpu\_id} là vị trí trong tệp gốc để tra ngược. Nhóm giữ
\PreColsSelected{} biến theo ý nghĩa đối với câu hỏi nghiên cứu, không theo tỉ lệ khuyết:
{\sloppy
\begin{itemize}
  \item biến mục tiêu \texttt{Max\_Power}; biến giải thích \texttt{TMUs} (số đơn vị xử lý
    texture, đại diện cho quy mô chip), \texttt{Memory}, \texttt{Memory\_Bandwidth},
    \texttt{Process};
  \item biến phụ \texttt{Core\_Speed} (xung nhịp), không bắt buộc phải có
    (mục~\ref{sec:missing});
  \item biến phân nhóm và định danh: \texttt{Name}, \texttt{Manufacturer},
    \texttt{Notebook\_GPU}, \texttt{Memory\_Type}, \texttt{Release\_Date},
    \texttt{Architecture};
  \item biến chỉ dùng khi xử lý: \texttt{Dedicated} (lọc GPU rời), \texttt{ROPs},
    \texttt{Texture\_Rate}, \texttt{Boost\_Clock}, \texttt{Memory\_Bus},
    \texttt{Memory\_Speed} (sửa và kiểm tra lỗi).
\end{itemize}
\par}
{\sloppy Các biến về cổng xuất hình, DirectX, giá bán\ldots{} không liên quan nên bị bỏ.
\texttt{PSU} và \texttt{Power\_Connector} cũng bị bỏ dù liên quan, vì nhà sản xuất chọn
chúng \emph{sau khi} đã biết công suất; đưa vào mô hình sẽ làm mô hình dự đoán tốt một
cách giả tạo. \texttt{Pixel\_Rate} trùng thông tin quy mô chip với \texttt{TMUs};
\texttt{Shader} là phiên bản shader, không phải số nhân xử lý.\par}

{\sloppy Trong tệp gốc, ô khuyết không để trống mà ghi chuỗi như \verb|"\n- "| hay
\verb|"\nUnknown Release Date "|; nếu đọc theo mặc định, \texttt{pandas} đếm được 0 ô
khuyết ở \texttt{Max\_Power}. Bước 3 xóa khoảng trắng ở hai đầu mỗi giá trị, đổi chuỗi
rỗng, \verb|"-"| và \verb|"Unknown Release Date"| thành \texttt{NA}
(\PreCellsMarkedMissing{} ô), và sửa \PreEncodingFixed{} tên card bị lỗi mã hóa
(\verb|"²"| thành \verb|"²"|). Vì vậy việc lọc ở bước 5 phải đặt sau bước làm sạch.\par}

\begin{lstlisting}[language=Python, caption={Đánh dấu ô khuyết (rút gọn từ hàm \texttt{clean\_text})}]
gpu = gpu.apply(lambda col: col.str.strip())   # "\n01-Mar-2009 " -> "01-Mar-2009"
gpu = gpu.replace(["", "-", "Unknown Release Date"], pd.NA)   # 0 -> 6474 NA cells
\end{lstlisting}

% ---------------------------------------------------------------------
\subsection{Chuyển định dạng và sửa lỗi (bước 4)}

{\sloppy Các biến số được ghi kèm đơn vị, như \verb|"141 Watts"| hay \verb|"64GB/sec"|. Mỗi biến
có một mẫu định dạng riêng để tách phần số; giá trị không khớp mẫu làm chương trình dừng,
nên không có giá trị nào bị đổi sai âm thầm. Cách xóa mọi ký tự không phải chữ số cho kết
quả sai, ví dụ \verb|"800MB/sec"| thành 800~GB/s thay vì 0,8~GB/s (\PreBandwidthInMb{}
dòng ghi theo MB/sec). \texttt{Release\_Date} được rút gọn thành năm
(\texttt{release\_year}); \texttt{Dedicated}, \texttt{Notebook\_GPU} thành đúng/sai;
\texttt{Manufacturer}, \texttt{Memory\_Type} thành biến phân loại. Hai lỗi dữ liệu được
xử lý theo nguyên tắc: giá trị không thể đúng thì không để nguyên, sửa nếu tìm được giá trị
đúng, nếu không thì bỏ \cite{vandenbroeck2005}.\par}

{\sloppy\textbf{Lỗi ``(x$N$)'' (sửa).} \PreRopsMultiplier{} dòng ghi \texttt{ROPs} dạng
``$a$ (x$N$)'', nghĩa là $a \times N$; ở đúng các dòng này \texttt{TMUs} và
\texttt{Memory\_Bus} bị ghi chia cho $N$. Ví dụ GTX~1080~Ti Zotac ghi \verb|"22 (x4)"| ROPs
và 56 TMUs, còn bản tham chiếu ghi 88 ROPs và 224 TMUs $= 56 \times 4$. Hai đẳng thức
kiểm chứng điều này (\texttt{check\_consistency}, cho lệch tối đa 10\%): tốc độ texture
$=$ TMUs $\times$ xung nhịp sai ở \PreTextureMismatchBefore{} trong \PreRopsMultiplier{}
dòng trước khi nhân lại (\PreTextureUntestable{} dòng không có xung nhịp để kiểm tra) và
\PreTextureMismatchAfter{} dòng sau đó; băng thông $=$ độ rộng bus $\times$ tốc độ bộ nhớ
$\times$ hệ số theo loại bộ nhớ sai ở cả \PreRopsMultiplier{} dòng trước khi sửa, còn sau
khi sửa cả tệp chỉ còn \PreBandwidthMismatchAfter{} dòng sai (cấu hình SLI, bị loại ở
bước 5). Nhóm nhân lại hai cột
với $N$; dữ liệu cuối còn \PreRepairedInFinal{} dòng đã sửa (năm \PreRepairedYears).
\texttt{Max\_Power} không bị ảnh hưởng: ở hai chip có cả dòng lỗi lẫn dòng không lỗi
(GP102 và Baffin PRO), công suất trung vị của hai loại dòng bằng nhau (220~W và 75~W).\par}

\begin{lstlisting}[language=Python, caption={Sửa lỗi ``(x$N$)'' (rút gọn từ hàm \texttt{parse\_columns})}]
# "22 (x4)" -> N = 4;  "88" -> N = 1
n = gpu["ROPs"].str.extract(r"\(x(\d+)\)")[0].astype(float).fillna(1)
gpu["tmus"] = pd.to_numeric(gpu["TMUs"]) * n
bus = gpu["Memory_Bus"].str.removesuffix(" Bit")      # "256 Bit" -> "256"
gpu["memory_bus_bit"] = pd.to_numeric(bus) * n
\end{lstlisting}

\textbf{Công suất 1~W (bỏ).} Radeon HD~6850 Club-3D OC Edition ghi \verb|"1 Watts"|,
trong khi \PreInvalidPowerSiblings{} card khác cùng chip Barts Pro đều ghi
\PreInvalidPowerSiblingWatts~W và tệp gốc ghi cho nó cùng bộ nguồn 500~W như bản tham
chiếu. Giá trị đúng không suy ra được từ chính dòng này nên ô được đổi thành \texttt{NA}
thay vì điền \PreInvalidPowerSiblingWatts~W. Nhóm không dùng ngưỡng chung kiểu ``dưới
10~W là lỗi'' vì \PreLowPowerLaptopCards{} card laptop ghi dưới 10~W và không có căn
cứ để coi chúng là lỗi. Ngoài ra, \PrePowerFarFromChipMedian{} card lệch hơn hai lần so với trung vị cùng chip được giữ
nguyên (có thể là thiết kế riêng của hãng gia công) và cần xem lại khi kiểm tra ngoại lai.

% ---------------------------------------------------------------------
\subsection{Lọc quan sát (bước 5)}

Đề tài nghiên cứu \emph{các card rời dùng một GPU, đã bán ra thị trường}, nên lần lượt
loại (Bảng~\ref{tab:rows}): GPU tích hợp (công suất không so sánh được với card rời) và
\PreUnknownDedicated{} dòng không rõ loại; cấu hình nhiều GPU (công suất ghi là của cả
hệ), nhận diện qua tên có ``SLI'', ``CrossFire'', ``X2''/``GX2'' sau số hiệu, dạng
``A + B'' hoặc thuộc danh sách đã biết (đã rà tay để không loại nhầm tên tản nhiệt như
``IceQ X2''); card chưa phát hành; dòng trùng hoàn toàn; và dòng thiếu biến bắt buộc, tức
mọi cột của dữ liệu cuối trừ xung nhịp.

\begin{table}[htbp]
  \centering
  \caption{Số dòng còn lại sau từng điều kiện lọc}
  \label{tab:rows}
  \small
  \begin{tabular}{l r r}
    \toprule
    \textbf{Điều kiện lọc} & \textbf{Số dòng} & \textbf{Bị loại} \\
    \midrule
    Dữ liệu gốc & \PreRowsRaw & \\
    Chỉ giữ GPU rời & \PreRowsDedicated
      & \the\numexpr\PreIntegrated+\PreUnknownDedicated\relax \\
    Loại cấu hình nhiều GPU & \PreRowsSingleGpu & \PreMultiGpuRemoved \\
    Loại card chưa phát hành & \PreRowsReleased & \PreUnreleasedRemoved \\
    Loại dòng trùng hoàn toàn & \PreRowsScoped & 0 \\
    Loại dòng thiếu biến bắt buộc & \PreRowsFinal & \PreIncompleteRemoved \\
    \bottomrule
  \end{tabular}
\end{table}

Dữ liệu cuối gồm \PreRowsFinal{} card (\PreDesktop{} desktop, \PreLaptop{} laptop), phát
hành từ \PreReleaseYearMin{} đến \PreReleaseYearMax, chỉ còn Nvidia, AMD và ATI. Đơn vị
quan sát là một sản phẩm card, không phải một chip: nhiều hãng gia công bán card chung chip
và thường ghi cùng công suất, nên các quan sát không hoàn toàn độc lập. Cần lưu ý điều
này khi diễn giải các kiểm định ở các mục sau.

% ---------------------------------------------------------------------
\subsection{Tạo biến mới và bộ dữ liệu cuối (bước 6)}

Bước này tạo các biến mà nhiều phân tích dùng chung một định nghĩa:
\texttt{segment} (Desktop/Laptop); \texttt{needs\_aux\_power} (công suất trên 75~W, mức
tối đa một khe PCIe cấp được, chỉ có ý nghĩa với card desktop);
\texttt{memory\_type\_group} (gộp các loại bộ nhớ dưới 30 card thành ``Other'', còn
\PreMemoryTypeGroups{} nhóm; nhóm ``Other'' chỉ \PreMemoryTypeOtherSize{} card, gồm các
công nghệ rất khác nhau nên chỉ mô tả, không so sánh); và
\texttt{texture\_rate\_per\_watt} (hiệu quả năng lượng, bằng tốc độ xử lý texture chia
cho công suất). Dữ liệu được lưu dưới dạng
\texttt{pandas.DataFrame}, mỗi dòng một card; Bảng~\ref{tab:columns} liệt kê các cột.

\begin{table}[htbp]
  \centering
  \caption{Các cột của bộ dữ liệu cuối và tỉ lệ khuyết trước khi loại dòng khuyết}
  \label{tab:columns}
  \small
  \begin{tabular}{l >{\ttfamily}l l r}
    \toprule
    \textbf{Vai trò} & \normalfont\textbf{Cột} & \textbf{Kiểu} & \textbf{Khuyết (\%)}$^{*}$ \\
    \midrule
    Định danh & gpu\_id, name & số nguyên, chuỗi & \PreMissingPctName \\
    Biến mục tiêu & max\_power\_w & số thực & \PreMissingPctMaxPowerW \\
    Tạo từ mục tiêu & needs\_aux\_power & đúng/sai & \PreMissingPctNeedsAuxPower \\
                    & texture\_rate\_per\_watt & số thực & \PreMissingPctTextureRatePerWatt \\
    Biến giải thích & tmus & số đếm & \PreMissingPctTmus \\
                    & memory\_bandwidth\_gb\_s & số thực & \PreMissingPctMemoryBandwidthGbS \\
                    & memory\_mb & số thực & \PreMissingPctMemoryMb \\
                    & process\_nm & số thực & \PreMissingPctProcessNm \\
    Biến phụ (được khuyết) & core\_speed\_mhz & số thực & \PreMissingPctCoreSpeedMhz \\
    Biến phân nhóm & manufacturer & phân loại & \PreMissingPctManufacturer \\
                   & segment & phân loại & \PreMissingPctSegment \\
                   & memory\_type(\_group) & phân loại & \PreMissingPctMemoryType \\
    Thời gian, chip & release\_year & số nguyên & \PreMissingPctReleaseYear \\
                    & architecture & chuỗi & \PreMissingPctArchitecture \\
    \bottomrule
  \end{tabular}

  \smallskip
  {\footnotesize $^{*}$ Trên \PreRowsScoped{} dòng thuộc phạm vi đề tài. Sau khi loại,
    chỉ \texttt{core\_speed\_mhz} còn khuyết (\PreCoreMissingInFinal{} dòng,
    \PreCoreMissingInFinalPct\%).\par}
\end{table}

% ---------------------------------------------------------------------
\subsection{Xử lý dữ liệu khuyết}
\label{sec:missing}

\textbf{Cách xử lý.} Nhóm không điền giá trị nào; dòng thiếu một biến bắt buộc bị loại;
riêng xung nhịp được phép khuyết. Kết quả là \PreRowsIncomplete{} dòng
(\PreRowsIncompletePct\%) bị loại, nhưng \PrePowerMissing{} dòng trong đó thiếu chính
công suất, tức biến mục tiêu; chỉ \PrePowerRowsLostByDropAll{} dòng có công suất mà bị
loại vì thiếu biến khác.

\textbf{Vì sao không điền.} Điền trung bình hay trung vị là tự tạo số liệu và làm kết
quả trông chắc chắn hơn thực tế: nếu điền trung vị cho công suất thì
\PrePowerAtMedianImputed\% số card có cùng một giá trị và độ lệch chuẩn giảm từ
\PrePowerSdDropped{} xuống \PrePowerSdImputed~W. Điền cho biến giải thích cũng làm sai
lệch mối liên hệ với công suất \cite{donders2006}. Tài liệu này cũng lưu ý việc loại dòng
có thể gây lệch khi ô khuyết không ngẫu nhiên và khuyến nghị điền nhiều lần (multiple
imputation); phương pháp đó vượt ngoài chương trình môn học, nên nhóm kiểm tra trực tiếp
mức lệch do loại dòng ở phần dưới.

\textbf{Vì sao xung nhịp không bắt buộc.} Nếu bắt buộc, dữ liệu chỉ còn
\PreRowsIfCoreRequired{} dòng, mất thêm \PreCoreMissingInFinal{} card đủ mọi thông số
khác; các card này có công suất thấp hơn (\PreCoreMissingPowerMean~W so với
\PreCoreKnownPowerMean~W) nên bỏ chúng làm dữ liệu lệch thêm về card công suất lớn. Xung
nhịp có liên hệ với công suất: trong cùng một chip, xung nhịp cao hơn 1\% đi kèm công suất
cao hơn khoảng \PreCoreWithinChipElasticity\% (ước lượng mô tả). Nhưng giữa các chip, xung
nhịp đi cùng tiến trình ($r = \PreCoreProcessCorr$) và năm phát hành
($r = \PreCoreYearCorr$); hai biến này giải thích khoảng \PreCoreExplainedPct\% sự khác
nhau về xung nhịp, nên khó tách riêng ảnh hưởng của nó. Vì vậy xung nhịp không nằm trong
mô hình chính; điều này không có nghĩa là xung nhịp không ảnh hưởng tới công suất, và hệ
số của tiến trình cần được hiểu là đã bao gồm phần đi kèm với xung nhịp. Đây là lựa chọn
đưa ra sau khi đã xem dữ liệu.

\textbf{Ô khuyết không ngẫu nhiên.} Card càng cũ càng hay thiếu công suất
(Bảng~\ref{tab:missing-period}): năm phát hành trung bình của card khuyết là
\PreMissingMeanYear{} so với \PreObservedMeanYear{} ($z = \PreMissingYearZ$,
$p$ \PreMissingYearP), và tỉ lệ AMD cao hơn (\PreMissingAmdPct\% so với
\PreObservedAmdPct\%, trong số card Nvidia và AMD). Các kiểm định $z$ coi mỗi card là độc lập nên chỉ dùng để mô tả.
Vì vậy dữ liệu cuối nghiêng về card đời mới: chỉ \PreRowsBeforeTwoThousandSix{} card trước
năm 2006, và các dòng có công suất nhưng bị loại có công suất trung bình
\PreLostPowerMean~W, thấp hơn \PreFinalPowerMean~W của dòng được giữ.

\begin{table}[htbp]
  \centering
  \caption{Tỉ lệ khuyết công suất và tỉ lệ bị loại theo năm phát hành
    (\PreRowsScoped{} dòng trước khi loại dòng khuyết)}
  \label{tab:missing-period}
  \small
  \begin{tabular}{l r r r}
    \toprule
    \textbf{Năm phát hành} & \textbf{Số card} & \textbf{Khuyết công suất (\%)}
      & \textbf{Bị loại (\%)} \\
    \midrule
    \csname @@input\endcsname generated/preprocessing_missing_by_period.tex
    Không rõ năm & \PreUnknownYearRows & \PreUnknownYearPowerMissingPct
      & \PreUnknownYearDroppedPct \\
    \bottomrule
  \end{tabular}
\end{table}


\textbf{Hệ quả.} Bộ dữ liệu cuối có công suất trung bình cao hơn một chút so với mọi card
có ghi công suất (\PreFinalPowerMean{} so với \PreAllPowerMean~W, Bảng~\ref{tab:power-groups}),
nên các kết luận ở mục sau chỉ đại diện cho GPU rời có đủ thông số, chủ yếu phát hành từ
2006 đến \PreReleaseYearMax. Chênh lệch desktop--laptop ổn định trên cả hai tập (desktop
gấp khoảng ba lần). Ngược lại, chênh lệch Nvidia trừ AMD đổi từ $\PreNvidiaMinusAmdFinal$
thành $\PreNvidiaMinusAmdAll$~W, nên kết luận về hãng cần được kiểm tra trên cả hai tập.

\begin{table}[htbp]
  \centering
  \caption{Công suất (W) trên bộ dữ liệu cuối và trên mọi card có ghi công suất}
  \label{tab:power-groups}
  \small
  \begin{tabular}{l r r r r r r}
    \toprule
    & \multicolumn{3}{c}{\textbf{Bộ dữ liệu cuối}}
      & \multicolumn{3}{c}{\textbf{Mọi card có công suất}} \\
    \cmidrule(lr){2-4} \cmidrule(lr){5-7}
    \textbf{Nhóm} & \textbf{Số card} & \textbf{TB} & \textbf{Trung vị}
      & \textbf{Số card} & \textbf{TB} & \textbf{Trung vị} \\
    \midrule
    Tất cả & \PreRowsFinal & \PreFinalPowerMean & \PreFinalPowerMedian
      & \PreRowsWithPower & \PreAllPowerMean & \PreAllPowerMedian \\
    \midrule
    \csname @@input\endcsname generated/preprocessing_power_by_group.tex
    \bottomrule
  \end{tabular}
\end{table}
